\section{Information Available to Each Agent}
\label{app:information}
A comparison across agent families is only fair if no agent learns more about the environment than the others. This appendix lists everything each agent receives.

\subsection{The interface}
Every agent, ENA and the baselines alike, interacts with an environment only through the Gymnasium interface. A reset returns an initial observation. A step returns the next observation, the reward, and whether the episode terminated or was truncated by the time limit. In CartPole the observation is the cart's position and velocity and the pole's angle and angular velocity. The physical parameters that distinguish the environments are never part of the observation. In CartPole these are gravity, the pole's mass and the pole's length; in HalfCheetah, gravity, friction and wind. No agent is told when the environment changes, which environment it is in, or how many environments there are. The environment names in our figures and tables come from the experiment's logs and are used only for analysis; ENA's niches are its own numbering.

\subsection{Task constants}
ENA reads two constants from the task definition: the episode time limit (\Res{cfg/cartpole/maxsteps} steps in CartPole) and the solved threshold ($R_{\text{solved}} = \Res{cfg/cartpole/solved}$). It uses them to rescale a truncated candidate return to a full-episode equivalent~\eqref{eq:plasticity} and as the utility scale of the specialist selection~\eqref{eq:bpr}. Both are public properties of the task, not of the physics; the baselines meet the same time limit through the truncation signal.

\subsection{Environment detection}
The detector's input is the sequence of transitions $(s_t, a_t, s_{t+1})$: the observation before acting, the action taken and the observation after, which is exactly what the agent sees. It does not read the reward, the physical parameters or the environment's name. A change of physics changes how the observed velocities evolve, so a niche's dynamics model~\eqref{eq:dynamics} stops predicting the observations within a few steps. This is why changes between the training environments are detected almost immediately and almost without error. Recognition is weaker in the cases this appendix makes explicit:
\begin{itemize}
\item \textit{Physics never seen in training.} At test the detector is frozen and creates no niches. When no niche explains the transitions within the margin $m$, the episode is left unrecognized and no niche is assigned. The selection~\eqref{eq:bpr} still weighs every niche by how well its dynamics model explains the episode, so it falls back on the closest known physics. Physics close to a known niche are recognized as that niche. Earth's and Neptune's gravity, for instance, share one.
\item \textit{False alarms.} In a stationary test phase every confirmed change is a false alarm. In physics near the margin $m$, the detector can alternate between recognizing a niche and leaving the physics unrecognized. Each confirmation resets the selection belief to uniform.
\item \textit{Fragmentation.} In some replicates one training environment is split over two niches.
\end{itemize}
Table~\ref{tab:detection} gives these numbers for the two ablation variants with environment detection. Every phase but the run's first follows a real switch, so at most one confirmed change per phase is a detection and the rest are false alarms. The first training phase, Mars, starts with no niche at all.

% GENERATED by app/src/tools/paper_results.py. Do not edit; re-run the script.
\begin{table}[htbp]
\caption{Environment Recognition in the Skill Ablation (30 Replicates)}
\begin{center}
\footnotesize
\setlength{\tabcolsep}{3pt}
\begin{adjustbox}{max width=\columnwidth}
\begin{tabular}{|l|l|l|c|c|c|c|}
\hline
\textbf{Agent} & \textbf{Phase} & \textbf{Env.} & \textbf{Unrec. (\%)} & \textbf{Changes} & \textbf{First (ep.)} & \textbf{Niches} \\
\hline
ENA-HoF+EnvDetection & Training & Mars & 0.4 $\pm$ 0.4 & 1.3 $\pm$ 1.9 & 243.0 $\pm$ 138.7 & 1.6 $\pm$ 0.6 \\
ENA-HoF+EnvDetection & Training & Jupiter & 0.2 $\pm$ 0.0 & 1.1 $\pm$ 0.3 & 1.0 $\pm$ 1.0 & 1.1 $\pm$ 0.3 \\
ENA-HoF+EnvDetection & Training & Earth & 0.1 $\pm$ 0.2 & 1.2 $\pm$ 0.4 & 0.5 $\pm$ 0.7 & 1.2 $\pm$ 0.4 \\
ENA-HoF+EnvDetection & Test & Mars & 9.8 $\pm$ 24.3 & 2.0 $\pm$ 10.6 & 6.3 $\pm$ 4.2 & -- \\
ENA-HoF+EnvDetection & Test & Neptune & 1.9 $\pm$ 5.7 & 2.4 $\pm$ 5.8 & 0.2 $\pm$ 0.5 & -- \\
ENA-HoF+EnvDetection & Test & Jupiter & 0.8 $\pm$ 0.8 & 1.0 $\pm$ 0.0 & 0.2 $\pm$ 0.4 & -- \\
ENA-HoF+EnvDetection & Test & Brown & 96.5 $\pm$ 11.4 & 6.8 $\pm$ 11.3 & 0.2 $\pm$ 0.5 & -- \\
ENA-HoF+EnvDetection & Test & Earth & 1.0 $\pm$ 0.9 & 0.0 $\pm$ 0.0 & -- & -- \\
\hline
ENA-Combined & Training & Mars & 0.3 $\pm$ 0.3 & 1.3 $\pm$ 2.5 & 202.4 $\pm$ 174.3 & 1.4 $\pm$ 0.5 \\
ENA-Combined & Training & Jupiter & 0.2 $\pm$ 0.1 & 1.1 $\pm$ 0.4 & 0.8 $\pm$ 0.9 & 1.1 $\pm$ 0.4 \\
ENA-Combined & Training & Earth & 0.0 $\pm$ 0.0 & 1.3 $\pm$ 0.5 & 0.0 $\pm$ 0.0 & 1.3 $\pm$ 0.5 \\
ENA-Combined & Test & Mars & 1.1 $\pm$ 1.1 & 0.8 $\pm$ 1.7 & 29.5 $\pm$ 68.8 & -- \\
ENA-Combined & Test & Neptune & 0.8 $\pm$ 1.1 & 1.2 $\pm$ 1.1 & 0.3 $\pm$ 0.5 & -- \\
ENA-Combined & Test & Jupiter & 0.6 $\pm$ 0.3 & 1.0 $\pm$ 0.0 & 0.2 $\pm$ 0.4 & -- \\
ENA-Combined & Test & Brown & 92.6 $\pm$ 12.6 & 16.4 $\pm$ 18.8 & 0.0 $\pm$ 0.0 & -- \\
ENA-Combined & Test & Earth & 0.4 $\pm$ 0.2 & 0.0 $\pm$ 0.0 & -- & -- \\
\hline
\end{tabular}
\end{adjustbox}
\label{tab:detection}
\end{center}
\end{table}


\subsection{Candidate evaluation}
ENA evaluates candidate specialists in fresh instances of the current environment, built by the same environment factory the experiment uses. It uses no physical parameter for any decision. Every candidate step is counted in ENA's interaction budget, so a candidate evaluation is equivalent to an extra episode in the current environment, with the same reset any episodic agent has. The baselines could spend the same budget in parallel copies of the environment, as vectorized training does. What ENA does assume is the ability to reset the environment and to run episodes that are not the one being played. That is free in simulation and holds for any episodic task, but it would cost real time on a physical system.

\subsection{Evaluation and test}
The first-solve checks (Section~\ref{sec:setup}) run on a separate copy of the current environment, for every agent alike, and their steps are not counted as training. At test, every agent is frozen: the baselines act greedily, and ENA's weights, its Plasticity ($\mathcal{P}=0$) and its dynamics models are fixed. The only thing that still adapts is the choice of gladiator. Every agent of a replicate meets the same seeded environments.

\section{Hyperparameters}
\label{app:hyperparameters}
\todo{Tables generated from the runner constants: ENA, DQN, PPO, SAC-Discrete, SAC; protocol constants.}

\section{Compute Environment}
\todo{Hardware, software versions (requirements.txt), seeds.}
